\documentclass{article}
\usepackage{amsmath}
\usepackage{amssymb}
\usepackage{amsthm}
\usepackage{mathtools}
\usepackage{geometry}
\usepackage{hyperref}

\newcommand{\rk}{\operatorname{rank}}
\newcommand{\GL}{\operatorname{GL}}
\newcommand{\R}{\mathbb{R}}
\newcommand{\Mat}{\mathrm{Mat}}
\newcommand{\tr}{\operatorname{tr}}

\theoremstyle{definition}
\newtheorem{theorem}{Theorem}[section]
\newtheorem{proposition}[theorem]{Proposition}
\newtheorem{lemma}[theorem]{Lemma}
\newtheorem{corollary}[theorem]{Corollary}
\newtheorem{definition}[theorem]{Definition}
\newtheorem{remark}[theorem]{Remark}
\newtheorem{example}[theorem]{Example}

\title{Gradient Equations in Weight Space}
\date{}

\begin{document}
\maketitle

\section{Direct Weight-Space Formulation}

We work directly on the parameter space
\[
\Omega = \prod_{\ell=1}^L \Mat_{d_\ell\times d_{\ell-1}}(\mathbb R),
\qquad
W=(W_1,\dots,W_L),
\]
with product map
\[
\mu(W)=W_L\cdots W_1 \in \Mat_{d_L\times d_0}(\mathbb R).
\]

For each layer $\ell$, we use the partial-product convention
\[
W_{>\ell}:=W_L\cdots W_{\ell+1}\in \Mat_{d_L\times d_\ell}(\mathbb R),
\qquad
W_{<\ell}:=W_{\ell-1}\cdots W_1\in \Mat_{d_{\ell-1}\times d_0}(\mathbb R),
\]
so that
\[
\mu(W)=W_{>\ell}W_\ell W_{<\ell},
\qquad
W_{>L}=I_{d_L},
\qquad
W_{<1}=I_{d_0}.
\]

\begin{definition}[Population quadratic loss]
The population loss is
\[
\mathcal L(W)
:=
\mathbb E \bigl[\|Y-\mu(W)X\|_2^2\bigr].
\]
Equivalently,
\[
\mathcal L(W)
=
\tr(\Sigma_Y)-2\tr(\mu(W)\Sigma_{XY})+\tr\bigl(\mu(W)\Sigma_X\mu(W)^\top\bigr).
\]
\end{definition}

\begin{proposition}[Gradient formula in weight space]
For each $\ell\in\{1,\dots,L\}$,
\[
\nabla_{W_\ell}\mathcal L(W)
=
2\,W_{>\ell}^{\top}\bigl(\mu(W)\Sigma_X-\Sigma_{YX}\bigr)W_{<\ell}^{\top}.
\]
\end{proposition}

\begin{proof}
The differential of the product map in the direction $\dot W_\ell$ is
\[
D\mu(W)[\dot W_\ell]
=
W_{>\ell}\dot W_\ell W_{<\ell}.
\]
Since
\[
\nabla_M\mathcal L(M)=2(M\Sigma_X-\Sigma_{YX})
\]
for the quadratic loss viewed as a function of the end-to-end map $M=\mu(W)$, the chain rule gives
\[
\nabla_{W_\ell}\mathcal L(W)
=
W_{>\ell}^{\top}\nabla_M\mathcal L(\mu(W))W_{<\ell}^{\top}
=
2\,W_{>\ell}^{\top}\bigl(\mu(W)\Sigma_X-\Sigma_{YX}\bigr)W_{<\ell}^{\top}.
\]
\end{proof}

\begin{proposition}[Transpose-gradient formula]
For each $\ell\in\{1,\dots,L\}$,
\[
\bigl(\nabla_{W_\ell}\mathcal L(W)\bigr)^\top
=
2\,W_{<\ell}\bigl(\Sigma_X\mu(W)^\top-\Sigma_{XY}\bigr)W_{>\ell}.
\]
\end{proposition}

\begin{proof}
Transpose the formula from the previous proposition and use the symmetry of $\Sigma_X$:
\[
\bigl(\mu(W)\Sigma_X-\Sigma_{YX}\bigr)^\top
=
\Sigma_X\mu(W)^\top-\Sigma_{XY}.
\]
\end{proof}

\begin{proposition}[Critical equations on the rank-$r$ sheet]
Assume that $W$ is a first-order critical point and let
\[
\mu(W)=P_S W^\ast
=
P_S\Sigma_{YX}\Sigma_X^{-1}
\]
for the unique subset $S$ given by the critical-point theorem. Then
\[
\Sigma_X\mu(W)^\top-\Sigma_{XY}
=
\Sigma_X\Sigma_X^{-1}\Sigma_{XY}P_S-\Sigma_{XY}
=
\Sigma_{XY}P_S-\Sigma_{XY}
=
-\Sigma_{XY}P_Q.
\]
Hence the transpose-gradient equations are equivalent to
\[
W_{<\ell}W_0 W_{>\ell}=0,
\qquad
1\leq \ell\leq L,
\]
where
\[
W_0:=W_S:=\Sigma_{XY}P_Q.
\]
\end{proposition}

\begin{remark}[Equivalent forms]
Since
\[
W^{\ast\top}=\Sigma_X^{-1}\Sigma_{XY},
\qquad
\tau_S:=W^{\ast\top}P_Q,
\qquad
W_0=\Sigma_X\tau_S,
\]
the same equations may also be written as
\[
W_{<\ell}\Sigma_XW^{\ast\top}P_QW_{>\ell}=0.
\]
If the data are whitened so that $\Sigma_X=I_{d_0}$, this simplifies to
\[
W_{<\ell}W^{\ast\top}P_QW_{>\ell}=0.
\]
\end{remark}

\begin{definition}[Weight-space cyclic equations]\label{def:weight-space-cyclic-equations}
For fixed $S$, define the matrices
\[
E_\ell(W):=W_{<\ell}W_0W_{>\ell}
\in \Mat_{d_{\ell-1}\times d_\ell}(\mathbb R).
\]
The cyclic equations in weight space are
\[
E_\ell(W)=0,
\qquad
1\leq \ell\leq L.
\]
\end{definition}

\begin{remark}
This is the direct weight-space replacement for the residual-space equations from the previous versions. The role of the fixed closing arrow is now played by the matrix
\[
W_0=W_S:=\Sigma_{XY}P_Q
\]
acting directly between the two endpoint spaces.
\end{remark}

\begin{example}[The case $L=3$]
When there are three weight matrices $W_1,W_2,W_3$, the definition gives exactly three vanishing relations:
\begin{align*}
E_1(W)=0 &\iff W_0W_3W_2 = 0, \\
E_2(W)=0 &\iff W_1W_0W_3 = 0, \\
E_3(W)=0 &\iff W_2W_1W_0 = 0.
\end{align*}

So the cyclic pattern is now visible directly on the matrices
\[
W_0,W_1,W_2,W_3
\]
with no extra transposes. The correct statement is not that all cyclic products
\[
W_kW_{k-1}W_{k-2}
\qquad (k \bmod 4)
\]
vanish. Rather, the relations are precisely the length-$L$ cyclic products that pass through the closing map $W_0$.

For $L=3$, the only length-$3$ cyclic product that is \emph{not} forced to vanish is
\[
W_3W_2W_1
=
\mu(W),
\]
which is consistent with Definition~\ref{def:weight-space-cyclic-equations} because the end-to-end map is not supposed to vanish on the rank-$r$ critical sheet.
\end{example}

\end{document}
